%%%%%%%%%%%%%%%%%%%%%%%%%%%%%%%%%%%%%%%%%%%%%%%%%%%%%%%%%%%%%
% ==================== 5.3：问题三的模型建立与求解 ====================
% ✓ 自查：前两问结果与本问条件衔接清楚；模型完整；步骤可复现；结论由真实结果支持。

\subsection{问题三的模型建立与求解}

\subsubsection{建模思路与模型选择}

问题三要求在前两问结果的基础上完成【概括问题三的任务】，需要同时考虑【关键条件一】与【关键条件二】。
本问将【问题一结果】和【问题二结果】分别作为【参数、输入或约束】，其核心是【关键数学问题】。
根据【问题结构、变量性质或求解规模】，选择【主模型名称】描述【目标与约束之间的关系】，
必要时以【基准方法或对比模型】检验所选方案的有效性。
若新增数据必须处理，则在此说明【处理方法及理由】，不固定设置数据预处理标题。

\subsubsection{模型的建立}

定义【本问新增变量】为【符号及含义】，并沿用前两问中的【已有变量或参数】。
依据【现实机理、数学原理或题目规则】建立【核心方程、目标函数或状态转移关系】，同时满足【主要约束】：

\begin{equation}
  \text{【问题三核心方程、目标函数或状态转移关系】},
  \label{eq:q3_core}
\end{equation}

\begin{equation}
  \text{【问题三约束、边界条件或适用条件】}.
  \label{eq:q3_constraint}
\end{equation}

\noindent\textbf{模型汇总：}
汇总前两问输入、本问目标和全部约束，得到问题三的完整模型

\begin{equation}
  \mathcal{M}_3:\quad
  \left\{
  \begin{aligned}
    &\text{【核心模型表达式】},\\
    &\text{【全部主要约束或演化规则】},\\
    &\text{【变量范围、初始条件或边界条件】}.
  \end{aligned}
  \right.
  \label{eq:q3_summary}
\end{equation}

\subsubsection{模型的求解}

\noindent\textbf{Step 1：}整合前两问输出与【本问新增条件】，形成统一的【参数集合、状态初值或候选方案集】。

\noindent\textbf{Step 2：}采用【算法或计算方法】依次完成【关键计算过程】，并记录【必要的中间结果】。

\noindent\textbf{Step 3：}依据【目标、终止条件或筛选准则】获得最终【预测、方案或机理参数】，并输出问题四所需的【结果】。

求解使用【软件、求解器或程序语言】完成，关键参数、初值和终止条件为【真实设置及依据】，完整程序见附录【编号】。

\subsubsection{求解结果与分析}

求解结果为【真实结果】。由【图或表编号】可见，【描述核心规律、方案特征或变化趋势】；
进一步比较【对象或情景】可知【差异及原因】。因此，问题三的结论为【直接回答问题三】，
其中【必要输出】将作为问题四综合分析的依据。

\subsubsection{模型检验与灵敏度分析}

根据【本问模型类型】，采用【检验方法】检查【关键数学关系、约束可行性或结果精度】。
真实检验结果【指标或现象】达到【判定标准或基准】，说明【检验结论】。
若【核心不确定参数】可能影响最终方案，则在【现实变化范围】内进行扰动，重点分析
【数值变化是否引起决策结构变化】，并报告【稳定区间、临界阈值或敏感方向】。
